\begin{abstract}
We present risk-based segmentation, a financial asset recommendation (FAR) approach that routes each customer into a per-band sub-model of LightGCN so that their declared Markets in Financial Instruments Directive II (MiFID II) risk band shapes recommendations during training. On top of this segmentation we ask where suitability is best enforced and compare two exclusive interventions: steering the model during training with a coherence margin loss, or making the data coherent in the first place by regrouping each customer onto the band implied by their purchases. This places risk suitability inside the model, whereas existing FAR systems optimise only return and ranking accuracy or correct for risk by post-hoc re-ranking. We evaluate on three axes, return on investment (ROI), ranking accuracy (nDCG), and risk suitability, which we make measurable with Profile Coherence at $k$ (PC@$k$), and summarise the three with their harmonic mean as a single balance score. An analytical study of FAR-Trans~\cite{sanzcruzado2024fartransinvestmentdatasetfinancial} motivates this axis. We find that 18.6\% of scoreable Buys lie two or more bands from the customer's declared tolerance, and this discordance is a persistent customer-level trait. A regression with cluster-robust standard errors then shows coherent Buys earn $+2.94$ percentage points higher 6-month realised return than discordant ones, so suitability is not a return tax. Across 69 temporal evaluation splits, the margin loss lifts PC@10 from 0.794 to 0.918 but trades away some accuracy, whereas regrouping lifts PC@10 to 0.949 while also improving nDCG and leaving ROI essentially unchanged, making it the only one of the two interventions to raise the overall balance. Risk suitability can therefore be optimised directly inside a recommender rather than bolted on afterwards, balanced against return and accuracy, offering a template for suitability-constrained recommendation in other regulated domains.
\end{abstract}
